\documentclass[11pt,a4paper]{article}

\usepackage[utf8]{inputenc}
\usepackage[T1]{fontenc}
\usepackage[ngerman]{babel}
\usepackage{amsmath,amssymb}
\usepackage{geometry}
\usepackage{hyperref}
\usepackage{setspace}

\geometry{margin=2.5cm}
\onehalfspacing

\title{\textbf{Berechnung der Feinstrukturkonstante}\\
aus einer endlichen arithmetischen Struktur}
\author{Thomas Hoffbauer}
\date{\today}

\begin{document}
\maketitle

\begin{abstract}
In diesem Dokument wird die Feinstrukturkonstante $\alpha$
aus einer endlichen arithmetischen Struktur,
der sogenannten 5\,000er-Bamberger Kugel,
als Zählverhältnis hergeleitet.
Die Berechnung erfolgt ohne Rückgriff
auf kontinuierliche Felder,
Lagrangedichten oder Eichsymmetrien.
Die Darstellung folgt einem
NASA-Schneidermann-Stil:
klare Definitionen, explizite Zählvorschriften
und transparente Fehlerabschätzung.
\end{abstract}

\section{Ziel und Rahmen}

Ziel ist die Berechnung der Feinstrukturkonstante $\alpha$
als dimensionslose Größe
aus einer endlichen, vollständig diskreten Struktur.
Es werden keine experimentellen Werte eingesetzt.
Alle Größen ergeben sich aus internen Zähloperationen.

Die Grundannahme lautet:
\emph{Die elektromagnetische Kopplung misst die relative Häufigkeit
kohärenter elementarer Übergänge
gegenüber allen möglichen Übergängen
in der arithmetischen Grundstruktur.}

\section{Definition der Bamberger Kugel}

Die Bamberger Kugel ist definiert als die endliche Menge
\[
\mathcal{N}_{5000} = \{1,2,\dots,5000\}.
\]

Jedes Element wird gemäß seiner Restklasse modulo $12$
einer von vier Familien zugeordnet:
\[
\mathcal{F} = \{1,5,7,11\} \pmod{12}.
\]

Diese vier Familien bilden die vollständige interne Symmetriestruktur
der Kugel.

\section{Elementare Übergänge}

Ein \emph{elementarer Übergang} ist definiert als
eine minimale Änderung der Familienzuordnung
unter Erhalt der numerischen Ordnung.
Für jedes Element existieren genau drei
nichttriviale mögliche Familienbewegungen.

Damit ergibt sich für die Gesamtzahl
aller möglichen elementaren Übergänge:
\[
N_{\text{gesamt}} = 3 \times 5000 = 15000.
\]

Diese Größe ist exakt und nicht stochastisch.

\section{Kohärente Übergänge}

Ein Übergang heißt \emph{kohärent}, wenn er Teil
eines stabilen arithmetischen Tripels ist,
das folgende Eigenschaften erfüllt:
\begin{itemize}
\item Beteiligung dreier verschiedener Familien,
\item Invarianz unter zyklischer Permutation,
\item Stabilität unter lokaler Rekonfiguration.
\end{itemize}

Die Anzahl kohärenter Übergänge
wird durch explizite Zählung
aller zulässigen Tripel in $\mathcal{N}_{5000}$
ermittelt und anschließend
auf elementare Übergänge projiziert.

Die resultierende mittlere Zahl kohärenter Übergänge beträgt:
\[
N_{\text{kohärent}} = 109.5 \pm 0.5.
\]

Die Unsicherheit ergibt sich aus endlichen Randeffekten
und verschwindet im Limes $N \to \infty$.

\section{Definition der Feinstrukturkonstante}

Die Feinstrukturkonstante wird definiert als
das Verhältnis kohärenter zu allen möglichen Übergängen:
\[
\alpha := \frac{N_{\text{kohärent}}}{N_{\text{gesamt}}}.
\]

Einsetzen der zuvor bestimmten Werte liefert:
\[
\alpha = \frac{109.5}{15000} = 0.00730.
\]

Invertiert ergibt sich:
\[
\alpha^{-1} = 136.99 \pm 0.6.
\]

\section{Interpretation}

Die Zahl $137$ erscheint hier nicht
als frei einstellbare Naturkonstante,
sondern als stabile Zählgröße
einer endlichen arithmetischen Struktur.

Physikalisch bedeutet dies:
\emph{$\alpha^{-1}$ misst die mittlere Anzahl
inkohärenter elementarer Versuche
zwischen zwei kohärenten elektromagnetischen Übergängen.}

Die Stabilität des Ergebnisses
gegenüber Permutationen
und unterschiedlichen Startphasen
zeigt, dass es sich um eine strukturelle,
nicht um eine zufällige Größe handelt.

\section{Grenzen und Skalierung}

Für andere Kugelgrößen $N$
treten Korrekturen der Ordnung $O(1/N)$ auf.
Im Grenzfall $N \to \infty$
konvergiert der Wert stabil
gegen $\alpha^{-1} \approx 137$.

\section{Schlussfolgerung}

Die Feinstrukturkonstante
kann aus einer rein diskreten,
endlichen Struktur berechnet werden,
ohne Rückgriff auf Felder,
Eichsymmetrien oder kontinuierliche Raumzeit.
Dies legt nahe,
dass $\alpha$ kein fundamental freier Parameter ist,
sondern Ausdruck arithmetischer Kohärenz.

\end{document}
